\chapter{Risk \& Capital Structure}\label{ch:risk-and-capital-structure}

% Scope: CAPM, beta, systematic vs idiosyncratic risk, Modigliani--Miller, WACC.
% Cross-link: WACC is the discount rate r in DCF (\cref{ch:valuation}) and in the CLV formula
%   (\cref{ch:customer-economics}).

\section{Systematic Risk \& CAPM}\label{sec:capm}
% complexity: 4 -> figure capm-sml

What return must a business earn to justify the risk its investors bear? CAPM answers by
decomposing risk into the part the market prices and the part it ignores, then deriving the
minimum required return accordingly. Investors can eliminate idiosyncratic (firm-specific)
risk for free via diversification. They cannot eliminate systematic risk --- the co-movement
with the broad market. Because only systematic risk survives in a diversified portfolio,
only systematic risk commands a return premium. Beta measures that co-movement; the security
market line prices it.

Let \(R_f\) be the risk-free rate, \(E(R_m)\) the expected market return, and \(\beta\)
the asset's systematic risk --- the slope of the regression of its excess returns on market
excess returns. The capital asset pricing model gives the required return:
\begin{equation}\label{eq:capm}
  E(R_i) \;=\; R_f \;+\; \beta\,\bigl[E(R_m)-R_f\bigr].
\end{equation}
The term \([E(R_m)-R_f]\) is the equity risk premium (ERP). For the US market over long
horizons the ERP is approximately \qty{5}{\percent} (geometric mean). Every asset with
\(\beta>1\) amplifies market moves; \(\beta<1\) dampens them.

CAPM rests on mean-variance preferences and normal return distributions. Empirical returns
are fat-tailed; Fama--French document that size and book-to-market explain return variation
that \(\beta\) misses. High-growth SaaS firms frequently have negative accounting earnings,
making historical beta estimates noisy (negative earnings produce unreliable return series).
Practitioners often use \emph{industry} beta --- unlever peers' betas, then re-lever at the
target capital structure --- to sidestep firm-specific noise.

\begin{example}[Deriving the cost of equity for the running SaaS]
The running B2B SaaS has estimated \(\beta=1.4\) (higher than the market, consistent with
the growth-stage volatility of software firms). The current risk-free rate (10-year
Treasury) is \qty{4}{\percent}; the ERP is \qty{5}{\percent}. Applying \cref{eq:capm}:
\[
  E(R_i) \;=\; \qty{4}{\percent} \;+\; 1.4\times\qty{5}{\percent}
          \;=\; \qty{4}{\percent} + \qty{7}{\percent}
          \;=\; \qty{11}{\percent}.
\]
This \qty{11}{\percent} is the cost of equity that feeds into WACC (\cref{eq:wacc})
and closes the book's discount-rate loop: it is the \(r\) in \cref{eq:gordon},
\cref{eq:pv}, \cref{eq:npv}, and \cref{eq:dcf} throughout the running example.
Managerial implication: raising growth guidance in a way that increases perceived
volatility --- lifting \(\beta\) from 1.4 to 1.7 --- adds \qty{1.5}{\percent} to
the required return, widening the denominator \((\mathrm{WACC}-g)\) from 0.08 to 0.095
and compressing the EV/Revenue multiple from roughly \(2.6\times\) to \(2.2\times\)
(\cref{eq:ev-revenue}).
\end{example}

\begin{figure}[htbp]
  \centering
  \includegraphics[width=0.86\linewidth]{capm-sml}
  \caption{The Security Market Line. Expected return rises linearly with \(\beta\);
  assets above the line are underpriced (positive alpha), assets below are overpriced.
  The running SaaS (\(\beta=1.4\)) plots at the \qty{11}{\percent} required return
  that anchors all DCF work in this book.}
  \label{fig:capm-sml}
\end{figure}

\begin{admonition}{Warning}
Treating \(\beta\) as a stable, intrinsic property. \emph{Beta is estimated from
historical co-movement}, which shifts with business model, leverage, and market
regime. A SaaS firm that was defensive in 2019 (\(\beta=0.8\)) may re-price to
\(\beta=1.6\) once interest rates rise and its growth premium becomes the
dominant risk factor. Refreshing beta at each capital-allocation cycle --- not
anchoring to a stale historical estimate --- is the discipline that keeps the
hurdle rate honest.
\end{admonition}

Multi-factor models (Fama--French three-factor, Carhart four-factor, AQR quality factor)
add size, value, and momentum loadings to CAPM's single \(\beta\). For private companies
without traded returns, analysts use a synthetic beta built from operating leverage (ratio
of fixed to total costs) and financial leverage --- both of which directly multiply the
systematic component of earnings volatility.

CAPM provides the equity-cost input that WACC needs; the two are inseparable.
\textcite{brealey2020} derive CAPM from mean-variance portfolio theory;
\textcite{ross2022} cover the empirical literature and factor alternatives;
\textcite{berk6thcorporate} is the standard reference for the industry-beta
methodology used in practice.

\section{The Weighted Average Cost of Capital}\label{sec:wacc}
% complexity: 3

The firm is funded by both equity and debt. Each has a different cost and tax treatment.
WACC is the single discount rate that correctly weights them for use in a DCF, an NPV
rule, or a CLV model. Equity holders bear residual risk and demand a return accordingly;
debt holders have a prior claim and accept a lower return, further subsidised by the tax
deductibility of interest. WACC blends these costs weighted by their share of total capital,
producing the true all-in cost of funding the firm's assets.

Let \(E\) and \(D\) be the market values of equity and debt, \(V=E+D\), \(r_e\) the cost
of equity from \cref{eq:capm}, \(r_d\) the pre-tax cost of debt, and \(T_c\) the corporate
tax rate. The tax shield on interest payments reduces the effective cost of debt from
\(r_d\) to \(r_d(1-T_c)\):
\begin{equation}\label{eq:wacc}
  \mathrm{WACC} \;=\; \frac{E}{V}\,r_e \;+\; \frac{D}{V}\,r_d\,(1-T_c).
\end{equation}
With \(r_e\) from \cref{eq:capm}, WACC becomes the unified discount rate that flows into
\cref{eq:dcf}, \cref{eq:terminal-value}, and \cref{eq:ev-revenue} in \cref{ch:valuation}.
It also equals \(r\) in the CLV formula \cref{eq:clv} in \cref{ch:customer-economics} and
in the growing perpetuity \cref{eq:gordon}: every present-value calculation in this book
uses the same underlying number.

WACC is stable only if the capital structure is stable. A firm executing a leveraged
recapitalisation, an LBO, or a staged debt paydown has a time-varying WACC; the adjusted
present value (APV) method --- which separately discounts the unlevered FCF and the PV of
the tax shield --- handles this more cleanly. Additionally, WACC treats the tax shield as
certain, which requires the firm to be profitable enough to use it. A growth-stage SaaS
with EBITDA near zero receives little tax-shield benefit, so the first equity-financed years
carry the full cost of equity.

\begin{example}[WACC with and without debt]
The equity-financed SaaS (\(D=0\)):
\[
  \mathrm{WACC} \;=\; \frac{E}{V}\times\qty{11}{\percent} \;+\; 0
               \;=\; \qty{11}{\percent}.
\]
Now suppose the founders consider raising \money{1000000} of debt at \qty{7}{\percent}
with a \qty{25}{\percent} tax rate while the firm is valued at \money{18000000}.
The new capital structure: \(E=\money{17000000}\), \(D=\money{1000000}\),
\(V=\money{18000000}\):
\[
  \mathrm{WACC} \;=\; \frac{17}{18}\times\qty{11}{\percent}
               \;+\; \frac{1}{18}\times\qty{7}{\percent}\times(1-0.25)
               \;\approx\; \qty{10.68}{\percent}.
\]
The reduction is real but small --- \qty{32}{\percent} of a percentage point ---
because the debt is modest relative to equity value. Modigliani--Miller
(\cref{sec:mm}) formalises why this near-irrelevance holds in its purest form,
and what breaks it in practice.
\end{example}

\begin{admonition}{Warning}
Using a \emph{book-value} capital structure instead of a market-value one in
\cref{eq:wacc}. Book values reflect historical accounting, not the opportunity
cost of today's capital. A firm with \money{5000000} of equity on the books but
a \money{50000000} market cap is predominantly equity-financed at market weights;
applying book weights dramatically overstates the equity fraction and overstates
WACC, causing the firm to reject positive-NPV projects.
\end{admonition}

For firms with complex capital structures (multiple tranches of debt, preferred equity,
warrants), each instrument enters \cref{eq:wacc} at its own cost and market weight.
Convertible notes require a split between debt and equity components (Black--Scholes value
of the embedded call) before weighting. Private-company WACC estimation adds another layer:
since neither equity nor debt is traded, both the cost of equity (using the industry-beta
method) and the cost of debt (using a synthetic rating from interest-coverage ratios) must
be estimated rather than observed.

WACC is the spine of the book's discount-rate architecture: its equity component comes from
\cref{eq:capm}; it feeds \cref{eq:dcf} and \cref{eq:terminal-value} in \cref{ch:valuation};
and it equals \(r\) in the customer-economics model (\cref{eq:clv}, \cref{eq:gordon}).
\textcite{brealey2020} derive the tax-shield adjustment; \textcite{modigliani1958capital}
is the original proposition.

\section{Modigliani--Miller}\label{sec:mm}
% complexity: 4

Does it matter how the firm is financed --- equity versus debt? M\&M provides the
conditions under which it does not, and equally, the imperfections that make capital
structure real. Understanding the theorem is a prerequisite for knowing when to ignore it.
In a world with no taxes, no bankruptcy costs, perfect information, and no agency frictions,
the total value of the firm is determined solely by its real assets --- not by how it slices
the claims on those assets into debt and equity. Financing is then a zero-sum
redistribution, not a value-creation act.

\textbf{Proposition I (no taxes):} \(V_L = V_U\) --- the value of a levered firm equals
the value of an identical unlevered firm. \textbf{Proposition II (no taxes):} as leverage
rises, equity holders demand a higher return to compensate for the increased risk borne:
\[
  r_e \;=\; r_A \;+\; \frac{D}{E}(r_A - r_d),
\]
where \(r_A\) is the unlevered (asset) cost of capital. With corporate taxes, the interest
tax shield has present value \(T_c\cdot D\), raising firm value:
\begin{equation}\label{eq:mm-wacc}
  V_L \;=\; V_U \;+\; T_c\,D,
  \qquad
  \mathrm{WACC} \;=\; r_A\,\bigl(1 - T_c\,D/V\bigr).
\end{equation}
\cref{eq:mm-wacc} shows WACC declining monotonically with leverage when taxes alone are
considered --- implying maximum value at \qty{100}{\percent} debt, which is evidently
wrong. The missing pieces are the subject of \cref{sec:optimal-capital-structure}.

M\&M's frictionless world omits three forces that dominate at moderate-to-high leverage:
(1) \emph{distress costs} --- legal, operational, and reputational costs triggered when
debt becomes hard to service; (2) \emph{agency costs} --- as debt rises, equity holders
gain an option-like incentive to take excessive risk (\textcite{jensen1976agency} formalise
this as the asset-substitution problem); (3) \emph{asymmetric information} --- the
pecking-order theory (\textcite{brealey2020}) predicts firms prefer retained earnings, then
debt, then equity, because issuing equity signals management believes the stock is
overvalued. Each friction drives a wedge between M\&M's irrelevance and observed financing
behaviour.

\begin{example}[M\&M tax shield for the SaaS]
At the SaaS's observed capital structure (\(D=0\), \(V=\money{18000000}\)) and
\qty{25}{\percent} tax rate, \cref{eq:mm-wacc} gives:
\[
  \mathrm{WACC} \;=\; \qty{11}{\percent}\times(1-0.25\times0) \;=\; \qty{11}{\percent}.
\]
The tax shield is zero because there is no debt. If the firm were to add
\money{5000000} of permanent debt:
\[
  \mathrm{WACC} \;\approx\; \qty{11}{\percent}\times\!\left(1-0.25\times\tfrac{5}{18}\right)
               \;\approx\; \qty{10.24}{\percent}.
\]
The lower WACC increases enterprise value mechanically, but only up to the point
where distress costs begin to offset the tax shield --- the trade-off that
\cref{sec:optimal-capital-structure} maps out.
\end{example}

\begin{admonition}{Warning}
Concluding from M\&M that capital structure \emph{never} matters. The theorem's
value is not its conclusion in the frictionless world but its enumeration of the
frictions that make capital structure consequential in the real world: taxes,
distress costs, agency costs, and asymmetric information. \emph{Every real-world
capital-structure decision is a bet on which friction dominates.}
\end{admonition}

The Jensen--Meckling agency-cost framework adds two further distortions beyond
\textcite{modigliani1958capital}: underinvestment (debt overhang reduces the equity
holder's incentive to fund positive-NPV projects because gains accrue to bondholders) and
overinvestment/risk-shifting (equity holders prefer risky projects because their position
resembles a call option with limited downside). These effects create an \emph{optimal} debt
level even absent taxes: the one that minimises total agency costs.

M\&M remains the foundation of capital-structure theory precisely because it forces clarity
about \emph{which} frictions matter in any given situation. \textcite{modigliani1958capital}
is the original; \textcite{jensen1976agency} extends to agency costs; \textcite{brealey2020}
provides the unified treatment for practitioners.

\section{Optimal Capital Structure}\label{sec:optimal-capital-structure}
% complexity: 5 -> figure capital-structure

How much debt should the firm carry? The trade-off theory says the answer is where the
marginal benefit of the tax shield equals the marginal cost of financial distress. The
pecking-order theory says the question is almost beside the point: firms finance
sequentially from internal cash first. Adding debt creates two opposing forces: the tax
shield raises firm value; distress costs lower it. The value-maximising leverage ratio sits
at the peak of the net benefit curve --- where the next dollar of debt adds more distress
cost than it saves in tax. For a profitable, stable firm this peak may sit at substantial
leverage; for a loss-making growth firm it sits near zero.

The trade-off theory defines optimal capital structure as:
\[
  V^* \;=\; V_U \;+\; PV(\text{tax shield}) \;-\; PV(\text{distress costs}).
\]
This is a qualitative framework: neither the distress-cost function nor the agency-cost
function has a closed-form solution, so the ``derivation'' is conceptual. \cref{eq:mm-wacc}
shows the tax shield rising monotonically with \(D/V\); distress costs are convex in
leverage (small at low debt, rising steeply once coverage ratios deteriorate). The optimum
is a level of \(D/V\) that cannot be improved by a marginal debt-for-equity swap.

Trade-off theory assumes the firm can predict its distress-cost function, which depends on
asset tangibility (hard assets are easier to liquidate, so costs are lower), earnings
stability, and the cost of renegotiation. High-growth intangible businesses (SaaS, biotech)
have high distress costs because their assets (customer relationships, IP, key personnel)
evaporate in distress --- pushing the optimal \(D/V\) close to zero. The pecking-order
theory offers an alternative explanation: firms do not optimise \(D/V\); they simply use
the cheapest available financing, which for profitable firms means internal cash first.

\begin{example}[Optimal structure for the growth-stage SaaS]
The growth-stage SaaS has negative EBITDA in years 1--2, no taxable income to shield, and
customer-relationship assets that would lose value rapidly in financial distress (churn
accelerates; talent departs). Trade-off theory therefore predicts an optimal
\(D/E\approx0\), which is precisely the observed financing: equity-financed via venture
capital, no debt. The \qty{11}{\percent} WACC is the pure cost of equity (\cref{eq:capm}).
Debt becomes relevant only once the firm reaches stable profitability and earns a tax shield
worth protecting. Managerial implication: the decision to remain unlevered is not timidity
--- it is the correct reading of the trade-off for a business whose value is concentrated
in intangibles and growth options.
\end{example}

\begin{figure}[htbp]
  \centering
  \includegraphics[width=0.86\linewidth]{capital-structure}
  \caption{Firm value versus leverage under the trade-off theory.
  The tax shield raises value above the unlevered baseline;
  distress costs (convex in leverage) erode it at high \(D/V\).
  The peak is the optimal capital structure. A growth-stage SaaS
  with intangible assets and volatile earnings sits near \(D/V=0\).}
  \label{fig:capital-structure}
\end{figure}

\begin{admonition}{Warning}
Benchmarking capital structure against industry averages without adjusting for
the firm's own stage and asset mix. A mature software firm with \qty{60}{\percent}
gross margins and stable free cash flow can support meaningful leverage;
a pre-profit SaaS cannot, even if it operates in the same industry. \emph{Industry
leverage ratios are a description, not a prescription.}
\end{admonition}

Market-timing theory (\textcite{brealey2020}) proposes that managers issue equity when
valuations are high and repurchase when they are low --- with no attention to an optimal
target ratio. The empirical evidence that leverage mean-reverts slowly (rather than quickly
toward a target) is consistent with both market-timing and adjustment costs. Dynamic
trade-off models allow the target \(D/V\) to shift as the firm moves from growth stage to
maturity, capturing the life-cycle pattern: VC equity → profitable debt capacity →
leveraged recapitalisation → buyback.

Capital structure is the bridge from the cost-of-capital framework developed in
\cref{sec:capm,sec:wacc} to the valuation mechanics of \cref{ch:valuation}: the right
capital structure minimises WACC and maximises firm value, which is what the DCF in
\cref{eq:dcf} and the multiples in \cref{eq:ev-revenue} are measuring.
\textcite{brealey2020} survey the empirical literature on capital-structure choice;
\textcite{jensen1976agency} develops the agency-cost rationale for an interior optimum.
